\documentclass{article}

\usepackage[a4paper,margin=1in]{geometry}

\usepackage{amsmath, amsthm, amssymb, amsfonts}
\usepackage{mathtools}
\usepackage{enumitem}
\usepackage{microtype}
\usepackage{xcolor}
\usepackage{fancyhdr}
\usepackage{graphicx}
\usepackage{hyperref}

% -------------------------------
% Fonts and Encoding
% -------------------------------
\usepackage{lmodern}
\usepackage[T1]{fontenc}
\usepackage[utf8]{inputenc}

% -------------------------------
% Theorem Environments
% -------------------------------
\theoremstyle{definition}
\newtheorem{definition}{Definition}[section]
\newtheorem{example}[definition]{Example}
\newtheorem{remark}[definition]{Remark}

\theoremstyle{plain}
\newtheorem{theorem}[definition]{Theorem}
\newtheorem{lemma}[definition]{Lemma}
\newtheorem{corollary}[definition]{Corollary}
\newtheorem{proposition}[definition]{Proposition}

% -------------------------------
% Header and Footer
% -------------------------------
\pagestyle{fancy}
\fancyhf{}
\rhead{\thepage}
\lhead{\textit{Mathematical Notes}}
\cfoot{}

% -------------------------------
% Custom Commands
% -------------------------------
\newcommand{\R}{\mathbb{R}}
\newcommand{\Z}{\mathbb{Z}}
\newcommand{\Q}{\mathbb{Q}}
\newcommand{\N}{\mathbb{N}}
\newcommand{\C}{\mathbb{C}}

\DeclareMathOperator{\Ker}{Ker}
\DeclareMathOperator{\Img}{Im}

% -------------------------------
% Examples & exercises
% -------------------------------
\newenvironment{solution}{\par\noindent\textbf{Solution.}\ }{\par}
\newenvironment{exercise}{\par\noindent\textbf{Exercise.}\ }{\par}

\begin{document}
	
	\title{Formulae}
	\author{Ramon A. Oba}
	\date{\today}
	
	\maketitle
	\section{A Structural Catalogue of Differentiation}
	
	Calculus begins with the study of local change.  
	The derivative is the object that measures this change, and its definition is grounded in the notion of limit.
	
	\subsection{The elemental definition of limit}
	
	Let \(f\) be a function defined on a punctured neighbourhood of a point \(a\).  
	We say that
	\[
	\lim_{x\to a} f(x)=L
	\]
	if for every \(\varepsilon>0\) there exists \(\delta>0\) such that
	\[
	0<|x-a|<\delta
	\quad \Longrightarrow \quad
	|f(x)-L|<\varepsilon.
	\]
	
	This means that \(f(x)\) can be made arbitrarily close to \(L\) by taking \(x\) sufficiently close to \(a\), without requiring \(x=a\).
	
	\subsection{The derivative as a limit}
	
	Let \(f\) be a function defined on an interval containing \(a\).  
	The derivative of \(f\) at \(a\) is defined by
	\[
	f'(a)
	=
	\lim_{h\to 0}\frac{f(a+h)-f(a)}{h},
	\]
	provided this limit exists.
	
	Equivalently,
	\[
	f'(a)
	=
	\lim_{x\to a}\frac{f(x)-f(a)}{x-a}.
	\]
	
	Thus the derivative is obtained by forming a difference quotient and passing to the limit as the input increment tends to zero.
	
	\subsection{Derivative table}
	
	The following formulas provide the basic stock of standard derivatives.
	
	\[
	\frac{d}{dx}(c)=0
	\]
	
	\[
	\frac{d}{dx}(x)=1
	\]
	
	\[
	\frac{d}{dx}(x^n)=nx^{n-1}
	\]
	
	\[
	\frac{d}{dx}\left(\frac{1}{x}\right)=-\frac{1}{x^2}
	\]
	
	\[
	\frac{d}{dx}(\sqrt{x})=\frac{1}{2\sqrt{x}}
	\]
	
	\[
	\frac{d}{dx}(e^x)=e^x
	\]
	
	\[
	\frac{d}{dx}(a^x)=a^x\ln a,
	\qquad a>0,\ a\neq 1
	\]
	
	\[
	\frac{d}{dx}(\ln x)=\frac{1}{x}
	\]
	
	\[
	\frac{d}{dx}(\sin x)=\cos x
	\]
	
	\[
	\frac{d}{dx}(\cos x)=-\sin x
	\]
	
	\[
	\frac{d}{dx}(\tan x)=\sec^2 x
	\]
	
	\[
	\frac{d}{dx}(\cot x)=-\csc^2 x
	\]
	
	\[
	\frac{d}{dx}(\sec x)=\sec x\tan x
	\]
	
	\[
	\frac{d}{dx}(\csc x)=-\csc x\cot x
	\]
	
	\[
	\frac{d}{dx}(\arctan x)=\frac{1}{1+x^2}
	\]
	
	\[
	\frac{d}{dx}(\arcsin x)=\frac{1}{\sqrt{1-x^2}}
	\]
	
	\[
	\frac{d}{dx}(\sinh x)=\cosh x
	\]
	
	\[
	\frac{d}{dx}(\cosh x)=\sinh x
	\]
	
	\[
	\frac{d}{dx}(\tanh x)=\operatorname{sech}^2 x
	\]
	
	These formulas form the primitive derivative layer.
	
	\subsection{Structural rules of differentiation}
	
	The subject becomes much cleaner when organized by the structural form of the expression being differentiated.
	
	\subsubsection{Level 1: Linearity}
	
	Differentiation is linear:
	\[
	\frac{d}{dx}\bigl(\alpha f(x)+\beta g(x)\bigr)
	=
	\alpha f'(x)+\beta g'(x).
	\]
	
	This handles sums and constant multiples.
	
	Example:
	\[
	\frac{d}{dx}(x^3+\sin x)=3x^2+\cos x.
	\]
	
	\subsubsection{Level 2: Affine inner map}
	
	A particularly important case is composition with an affine map
	\[
	g(x)=ax+b.
	\]
	
	Let \(F\) be differentiable. Then
	\[
	\frac{d}{dx}(F\circ g)(x)
	=
	F'(g(x))g'(x)
	=
	a\,F'(ax+b).
	\]
	
	Equivalently,
	\[
	\boxed{
		\frac{d}{dx}F(ax+b)=a\,F'(ax+b)
	}
	\]
	
	This yields a clean and unified catalogue for derivatives of affine compositions.
	
	Examples:
	\[
	\frac{d}{dx}\sin(2x)=2\cos(2x)
	\]
	
	\[
	\frac{d}{dx}e^{5x-1}=5e^{5x-1}
	\]
	
	\[
	\frac{d}{dx}\ln(3x+7)=\frac{3}{3x+7}
	\]
	
	\[
	\frac{d}{dx}(4x-1)^n
	=
	4n(4x-1)^{n-1}
	\]
	
	\subsubsection{Level 3: General composition}
	
	More generally, if \(g\) is differentiable, then
	\[
	\boxed{
		\frac{d}{dx}F(g(x))=F'(g(x))g'(x)
	}
	\]
	
	This is the general chain rule.
	
	Example:
	\[
	\frac{d}{dx}\sin(x^2)=\cos(x^2)\cdot 2x.
	\]
	
	Thus the affine case is a special case of the general chain rule.
	
	\subsubsection{Product rule}
	
	When the expression has product structure,
	\[
	\boxed{
		\frac{d}{dx}\bigl(u(x)v(x)\bigr)=u'(x)v(x)+u(x)v'(x)
	}
	\]
	
	Example:
	\[
	\frac{d}{dx}(xe^x)=e^x+xe^x.
	\]
	
	\subsubsection{Quotient rule}
	
	When the expression has quotient structure,
	\[
	\boxed{
		\frac{d}{dx}\left(\frac{u(x)}{v(x)}\right)
		=
		\frac{u'(x)v(x)-u(x)v'(x)}{v(x)^2}
	}
	\qquad v(x)\neq 0.
	\]
	
	Example:
	\[
	\frac{d}{dx}\left(\frac{x}{x^2+1}\right)
	=
	\frac{(1)(x^2+1)-x(2x)}{(x^2+1)^2}
	=
	\frac{1-x^2}{(x^2+1)^2}.
	\]
	
	\subsection{Conceptual summary}
	
	Differentiation is organized by a small number of structural principles.  
	The main task is to identify the form of the expression and then apply the simplest compatible rule.
	
	\begin{center}
		\begin{tabular}{ll}
			\textbf{Structure of the expression} & \textbf{Differentiation rule} \\[4pt]
			sum & linearity \\[4pt]
			composition with \(ax+b\) & affine chain rule \\[4pt]
			general composition & chain rule \\[4pt]
			product & product rule \\[4pt]
			quotient & quotient rule
		\end{tabular}
	\end{center}
	
	\vspace{1em}
	
	\section{A Structural Catalogue of Indefinite Integration}
	
	Indefinite integration studies the inverse problem: given a function \(f\), find a function \(F\) such that
	\[
	F'(x)=f(x).
	\]
	Such a function \(F\) is called an antiderivative of \(f\), and one writes
	\[
	\int f(x)\,dx=F(x)+C,
	\]
	where \(C\) is an arbitrary constant.
	
	The cleanest way to organize integration is to begin with a stock of known antiderivatives and then see how these antiderivatives are transported through the basic structural forms of the integrand.
	
	\subsection{Antiderivative table}
	
	The following formulas provide the basic stock of standard antiderivatives.
	
	\[
	\int 1\,dx = x + C
	\]
	
	\[
	\int x^n\,dx = \frac{x^{n+1}}{n+1}+C,\qquad n\neq -1
	\]
	
	\[
	\int \frac{1}{x}\,dx = \ln|x|+C
	\]
	
	\[
	\int e^x\,dx = e^x+C
	\]
	
	\[
	\int a^x\,dx = \frac{a^x}{\ln a}+C,
	\qquad a>0,\ a\neq 1
	\]
	
	\[
	\int \cos x\,dx = \sin x+C
	\]
	
	\[
	\int \sin x\,dx = -\cos x+C
	\]
	
	\[
	\int \sec^2 x\,dx = \tan x+C
	\]
	
	\[
	\int \csc^2 x\,dx = -\cot x+C
	\]
	
	\[
	\int \sec x\tan x\,dx = \sec x+C
	\]
	
	\[
	\int \csc x\cot x\,dx = -\csc x+C
	\]
	
	\[
	\int \frac{1}{1+x^2}\,dx = \arctan x+C
	\]
	
	\[
	\int \frac{1}{\sqrt{1-x^2}}\,dx = \arcsin x+C
	\]
	
	\[
	\int \cosh x\,dx = \sinh x+C
	\]
	
	\[
	\int \sinh x\,dx = \cosh x+C
	\]
	
	\[
	\int \operatorname{sech}^2 x\,dx = \tanh x+C
	\]
	
	These formulas form the primitive integral layer.
	
	\subsection{Structural rules of indefinite integration}
	
	The principal structures are:
	\[
	\text{sum},\qquad \text{composition},\qquad \text{product},\qquad \text{rational quotient}.
	\]
	These correspond respectively to:
	\[
	\text{linearity},\qquad \text{substitution},\qquad \text{integration by parts},\qquad \text{algebraic decomposition}.
	\]
	
	\subsubsection{Level 1: Linearity}
	
	Since differentiation is linear, integration is linear as well:
	\[
	\int \bigl(\alpha f(x)+\beta g(x)\bigr)\,dx
	=
	\alpha\int f(x)\,dx+\beta\int g(x)\,dx.
	\]
	
	Example:
	\[
	\int (x^2+\cos x)\,dx
	=
	\frac{x^3}{3}+\sin x+C.
	\]
	
	\subsubsection{Level 2: Affine inner map}
	
	A particularly important case is composition with an affine map
	\[
	g(x)=ax+b,
	\qquad a\neq 0.
	\]
	
	Let \(f\) be a function with antiderivative \(F\), so that
	\[
	F'(u)=f(u).
	\]
	Then
	\[
	\frac{d}{dx}(F\circ g)(x)
	=
	F'(g(x))g'(x)
	=
	a\,f(g(x)).
	\]
	Hence
	\[
	\int (f\circ g)(x)\,dx
	=
	\frac{1}{a}(F\circ g)(x)+C.
	\]
	
	Equivalently,
	\[
	\boxed{
		\int f(ax+b)\,dx
		=
		\frac{1}{a}F(ax+b)+C
	}
	\qquad (a\neq 0).
	\]
	
	Examples:
	\[
	\int \sin(2x)\,dx
	=
	-\frac{1}{2}\cos(2x)+C
	\]
	
	\[
	\int e^{5x-1}\,dx
	=
	\frac{1}{5}e^{5x-1}+C
	\]
	
	\[
	\int \frac{1}{3x+7}\,dx
	=
	\frac{1}{3}\ln|3x+7|+C
	\]
	
	\[
	\int (4x-1)^n\,dx
	=
	\frac{(4x-1)^{n+1}}{4(n+1)}+C,
	\qquad n\neq -1
	\]
	
	\subsubsection{Level 3: General substitution}
	
	If \(F'=f\), then
	\[
	\frac{d}{dx}F(g(x))
	=
	f(g(x))g'(x).
	\]
	Therefore
	\[
	\boxed{
		\int f(g(x))g'(x)\,dx
		=
		F(g(x))+C
	}
	\]
	
	Example:
	\[
	\int \cos(x^2)\,2x\,dx
	=
	\sin(x^2)+C.
	\]
	
	Thus the affine rule is a special case of general substitution.
	
	\subsubsection{Integration by parts}
	
	When the integrand has product structure, one uses the product rule in reverse:
	\[
	(uv)'=u'v+uv'.
	\]
	Hence
	\[
	\boxed{
		\int u\,dv = uv-\int v\,du
	}
	\]
	
	Example:
	\[
	\int xe^x\,dx.
	\]
	Take
	\[
	u=x,\qquad dv=e^x\,dx.
	\]
	Then
	\[
	du=dx,\qquad v=e^x,
	\]
	so
	\[
	\int xe^x\,dx
	=
	xe^x-\int e^x\,dx
	=
	xe^x-e^x+C.
	\]
	
	Another example is
	\[
	\int \ln x\,dx
	=
	x\ln x-x+C.
	\]
	
	\subsubsection{Rational fractions and partial fractions}
	
	When the integrand is a rational function
	\[
	\frac{P(x)}{Q(x)},
	\]
	with \(P\) and \(Q\) polynomials, the method is primarily algebraic.
	
	\paragraph{Step 1: Polynomial division.}
	If
	\[
	\deg P \geq \deg Q,
	\]
	perform polynomial division:
	\[
	\frac{P(x)}{Q(x)}
	=
	S(x)+\frac{R(x)}{Q(x)},
	\qquad \deg R<\deg Q.
	\]
	
	\paragraph{Step 2: Factor the denominator.}
	Factor \(Q(x)\) as far as possible.
	
	\paragraph{Step 3: Decompose into partial fractions.}
	Write the proper rational expression as a sum of simpler terms.
	
	\paragraph{Step 4: Integrate term by term.}
	Use the antiderivative table together with linearity and substitution when needed.
	
	If the denominator has distinct linear factors,
	\[
	\frac{P(x)}{(x-a)(x-b)},
	\]
	use
	\[
	\frac{P(x)}{(x-a)(x-b)}
	=
	\frac{A}{x-a}+\frac{B}{x-b}.
	\]
	
	If the denominator has a repeated linear factor,
	\[
	\frac{P(x)}{(x-a)^n},
	\]
	use
	\[
	\frac{A_1}{x-a}+\frac{A_2}{(x-a)^2}+\cdots+\frac{A_n}{(x-a)^n}.
	\]
	
	If the denominator contains an irreducible quadratic factor,
	\[
	x^2+px+q,
	\]
	use
	\[
	\frac{Ax+B}{x^2+px+q}.
	\]
	
	Example:
	\[
	\int \frac{1}{x^2-1}\,dx.
	\]
	Since
	\[
	x^2-1=(x-1)(x+1),
	\]
	write
	\[
	\frac{1}{x^2-1}
	=
	\frac{1}{2(x-1)}-\frac{1}{2(x+1)}.
	\]
	Hence
	\[
	\int \frac{1}{x^2-1}\,dx
	=
	\frac{1}{2}\ln|x-1|
	-
	\frac{1}{2}\ln|x+1|
	+C.
	\]
	
	Another example:
	\[
	\int \frac{x}{x^2+1}\,dx.
	\]
	Since the numerator is proportional to the derivative of the denominator,
	\[
	u=x^2+1,\qquad du=2x\,dx,
	\]
	and therefore
	\[
	\int \frac{x}{x^2+1}\,dx
	=
	\frac{1}{2}\ln(x^2+1)+C.
	\]
	
	\subsection{Conceptual summary}
	
	Indefinite integration is not a disconnected collection of tricks.  
	It is governed by a small number of structural principles, each corresponding to a clear algebraic form of the integrand.
	
	\begin{center}
		\begin{tabular}{ll}
			\textbf{Structure of the integrand} & \textbf{Integration method} \\[4pt]
			sum & linearity \\[4pt]
			composition with \(ax+b\) & affine substitution \\[4pt]
			general composition \(f(g(x))g'(x)\) & substitution \\[4pt]
			product & integration by parts \\[4pt]
			rational quotient & algebraic decomposition / partial fractions
		\end{tabular}
	\end{center}
	
	\vspace{1em}
	
	\section{Shared Structural Table}
	
	The two catalogues are best seen together.
	
	\begin{center}
		\begin{tabular}{lll}
			\textbf{Structure} & \textbf{Differentiation} & \textbf{Integration} \\[4pt]
			sum & linearity & linearity \\[4pt]
			composition with \(ax+b\) & affine chain rule & affine substitution \\[4pt]
			general composition & chain rule & substitution \\[4pt]
			product & product rule & integration by parts \\[4pt]
			quotient / rational form & quotient rule & algebraic decomposition / partial fractions
		\end{tabular}
	\end{center}
	
	Thus both subjects become much cleaner when organized structurally: one identifies the algebraic form first, and only then applies the corresponding rule.
	
\end{document}